\section{Guarantees}\label{sec:theory}

All results are finite-sample and hold for any fixed score function, so they do not depend on the quality of the base model or of the predictor; proofs are in Appendix~\ref{app:proofs}. Each concerns the expected coverage of a random variant of a new family, not the coverage of every family. Table~\ref{tab:guarantees} lists what each one does and does not cover.

\begin{proposition}[Family-level calibration]\label{prop:hcp}
Under A1, the set~\eqref{eq:threshold} with $\omega\equiv1$ satisfies $\Cov(C)\ge1-\alpha$.
\end{proposition}
Proposition~\ref{prop:hcp} is Theorem~1 of~\cite{lee2026hierarchical} with the test point replaced by a uniformly random variant of the test family, which removes the need for exchangeability within families.

\begin{theorem}[Coverage over a class of family-level shifts]\label{thm:conditional}
Let $P$ be a law of families and $\Phi$ fixed. For $f\in\mathcal{F}$ with $f(D)\ge0$ $P$-almost surely and $0<\E_Pf(D)<\infty$, let $Q_f$ be the family law with $dQ_f/dP=f(D)/\E_Pf(D)$. If $G_1,\dots,G_K$ are i.i.d.\ $P$ and, independently, $G_{K+1}\sim Q_f$, the set~\eqref{eq:condcal} satisfies $\Cov(C)\ge1-\alpha$. The same holds under A1 ($f\equiv1$; the families need only be exchangeable).
\end{theorem}

\begin{corollary}[Any family-level shift]\label{cor:conditional}
With $G_1,\dots,G_K$ i.i.d.\ $P$ as in Theorem~\ref{thm:conditional} and a test family with any law $Q$, $\Cov(C)\ge1-\alpha-\inf_f\dtv(Q,Q_f)$, the infimum over the $f$ of Theorem~\ref{thm:conditional}. If $\Phi=(1,\hat w)$ for a ratio estimate $\hat w\ge0$ independent of the families, with $0<\E_P\hat w<\infty$ and $d\hat Q/dP=\hat w(D)/\E_P\hat w(D)$, then
\[
\E\big[\gamma_{K+1}(C)\,\big|\,\hat w\big]\;\ge\;1-\alpha-\min\big\{\dtv(Q,P),\ \dtv(Q,\hat Q)\big\}.
\]
\end{corollary}
Here $\dtv(Q,P)$ bounds the loss of family-level calibration alone (the same argument applied to Proposition~\ref{prop:hcp}) and $\dtv(Q,\hat Q)$ that of reweighting by $\hat w$ (Proposition~\ref{prop:weighted}), so conditioning on the estimated ratio is never worse, in these bounds, than either. Both distances are unobservable, so the corollary motivates conditioning on both but certifies nothing under an arbitrary shift. Conditioning needs no total-variation term under family exchangeability, when $\hat w$ is the true ratio up to scale, and for every shift whose ratio is a nonnegative function in the span of $1$ and $\hat w$; with descriptor features (ProtCOVER-D), for every shift whose ratio is nonnegative and affine in the descriptors, without target families.

\emph{What is new.} Table~\ref{tab:guarantees} places these results among related guarantees. A family's miss rate $1-\hat F_k(q)$ is a bounded loss that does not increase with the threshold, and family-weighted coverage is one minus its risk. HCP's threshold is therefore conformal risk control~\cite{angelopoulos2024conformal} with families as units, and~\eqref{eq:condcal} is the conditional risk control of Blot et al.~\cite{blot2025automatically}, which carries the conditional calibration of Gibbs et al.~\cite{gibbs2025conformal} to such losses, applied to families with the test family's loss imputed as one; constant features give HCP's threshold (up to ties; Appendix~\ref{app:proofs}). Pooling families is thus the risk-control view applied to grouped data, not a new device. What Theorem~\ref{thm:conditional} adds is that the program is valid as computed, where the test family's labels are unknown, its scores sit at $+\infty$, the minimizer need not be unique and the program can be unbounded: the largest optimal fitted value dominates the threshold of the program with the test family's true scores (step (d) of the proof). This adapts Gibbs et al.'s Proposition~6, which imputes an upper bound for the unknown scores and assumes min-norm solutions; Blot et al.\ treat neither ties nor unbounded programs, and Section~\ref{sec:method} computes the threshold exactly. Conditioning on the span of a constant and a cross-fitted ratio (Gibbs et al.\ note a single estimated shift as a one-element class) then gives Corollary~\ref{cor:conditional}. Applying Gibbs et al.\ to one random variant per family, the analogue of subsampling once~\cite{dunn2023distribution}, also gives family-weighted coverage over $\mathcal{F}$, because the subsampled pairs are exchangeable points, but discards all but one calibration variant of each family; we use it as a baseline. Duchi et al.~\cite[Cor.~7.27]{duchi2024predictive} regress one coverage threshold per environment on environment features, which gives a new family the conservative level $(1-\alpha)(1-\delta)$, without pooling its variants.

\begin{proposition}[Reweighting variant]\label{prop:weighted}
Under A2, the set~\eqref{eq:threshold} with $\omega_k=w(D_k)$, $k\le K+1$, satisfies $\Cov(C)\ge1-\alpha$. With $\omega_k=\hat w(D_k)$ for $\hat w\ge0$ independent of the families with $0<\E_P\hat w(D)<\infty$ and $\hat Q$ as in Corollary~\ref{cor:conditional}, $\E[\gamma_{K+1}(C)\mid\hat w]\ge1-\alpha-\dtv(Q,\hat Q)$.
\end{proposition}
Proposition~\ref{prop:weighted} lifts weighted exchangeability~\cite{tibshirani2019conformal} and its estimated-weight bound~\cite[Thm.~3]{lei2021conformal} from variants to families.

\begin{lemma}[Cross-fitting]\label{lem:crossfit}
For each fold $j$, the ratio $\hat w_j$ is a function only of the fitting calibration families and of the target families outside fold $j$, standardization, rescaling and clipping included. With independent families (as under A2) it is therefore independent of the calibration families used in~\eqref{eq:threshold} and~\eqref{eq:condcal} and of every family in fold $j$, and Corollary~\ref{cor:conditional} and Proposition~\ref{prop:weighted} hold for the queries of fold $j$ conditionally on $\hat w_j$. If instead all non-training families are exchangeable, as under a random split, then conditionally on the families that fit $\hat w_j$ the remaining calibration families and each family of fold $j$ are exchangeable, and Theorem~\ref{thm:conditional} with $f\equiv1$ holds conditionally on $\hat w_j$.
\end{lemma}

\begin{proposition}[Class-conditional versions]\label{prop:mondrian}
For $y\in\{0,1\}$ such that a family has $N_{y}>0$ with positive probability under $P$ and under the test law, the Mondrian versions of Propositions~\ref{prop:hcp} and~\ref{prop:weighted}, Theorem~\ref{thm:conditional}, Corollary~\ref{cor:conditional} and Lemma~\ref{lem:crossfit} hold with $\Cov_y$ in place of $\Cov$ and with $Q$, $P$, $Q_f$, $\hat Q$ conditioned on $N_{K+1,y}>0$.
\end{proposition}

\begin{lemma}[Enlargement]\label{lem:enlarge}
If $C'(x)\supseteq C(x)$ for all $x$ almost surely, where $C'$ may depend on all observed data, then $\Cov(C')\ge\Cov(C)$ and $\Cov_y(C')\ge\Cov_y(C)$.
\end{lemma}
The gate (iii) only replaces a set by $\{0,1\}$, so by Lemma~\ref{lem:enlarge} it preserves every guarantee above, whatever its parameters; it gives no guarantee for the families it retains.

\begin{table}[t]
\caption{Guarantees of related methods and of ProtCOVER's versions: finite-sample lower bounds of $1-\alpha$ on the target under the stated conditions, unless noted (repeated subsampling: $1-2\alpha$; multi-environment: with probability $1-\delta$); ``bound'': reduced by a total-variation distance that unlabeled data cannot estimate. ProtCOVER's versions hold for $\Cov$ and for $\Cov_y$ given $N_{K+1,y}>0$ (Proposition~\ref{prop:mondrian}) and, with cross-fitted ratios, conditionally on each fold's ratio (Lemma~\ref{lem:crossfit}).}
\label{tab:guarantees}
\centering\scriptsize\setlength{\tabcolsep}{3pt}
\begin{tabular}{p{3.3cm}lp{2.5cm}p{5.3cm}p{1.9cm}}
\toprule
Method & Unit & Target & Guaranteed under & Target data \\
\midrule
Split, Mondrian~\cite{lei2018distribution,vovk2012conditional} & variant & a new variant (of class $y$) & exchangeable variants & none \\
Weighted~\cite{tibshirani2019conformal,lei2021conformal} & variant & a new variant & covariate shift with the true ratio; estimated: bound & unlabeled \\
GWCP~\cite{bhattacharyya2024group} & variant & a new variant & shift of the proportions of known groups & group proportions \\
Conditional calibration~\cite{gibbs2025conformal,blot2025automatically} & variant & a new variant & every shift with ratio $f(X)$, $f\in\mathcal{F}$, $f\ge0$ & none \\
HCP, subsampling~\cite{lee2026hierarchical,dunn2023distribution} & family & a variant of a new family & A1 & none \\
Multi-environment~\cite[Cor.~7.27]{duchi2024predictive} & environment & a new environment, w.p.\ $1-\delta$ & shifts in a class of environment features & none \\
Conditional calibration, one variant per family & family & $\Cov$ & every shift with ratio $f(D)$, $f\in\mathcal{F}$, $f\ge0$ (incl.\ A1) & as $\mathcal{F}$ needs \\
\midrule
ProtCOVER-W (Prop.~\ref{prop:weighted}) & family & $\Cov$, $\Cov_y$ & A2 with the true ratio; estimated: bound & unlabeled families \\
ProtCOVER-D (Thm.~\ref{thm:conditional}) & family & $\Cov$, $\Cov_y$ & every shift with ratio affine in the descriptors, $\ge0$ (incl.\ A1) & none \\
ProtCOVER (Thm.~\ref{thm:conditional}, Cor.~\ref{cor:conditional}) & family & $\Cov$, $\Cov_y$ & every shift with ratio $a+b\,\hat w(D)\ge0$ (incl.\ A1); otherwise bound & unlabeled families \\
Gate (Lemma~\ref{lem:enlarge}) & family & & never lowers coverage; nothing on the retained families & \\
\midrule
Not guaranteed & \multicolumn{4}{p{10.6cm}}{coverage of any given family; the splits S2--S5, where A2 holds at best approximately and is violated in S5} \\
\bottomrule
\end{tabular}
\end{table}
